\section{Finite computations and split expressions}\label{alg:section}

The language problem will be represented by a finite computation. This
section supplies three interfaces for manipulating it: an expression
identity that propagates an update through episodes, a Boolean operation
that removes an affine correction, and a finite table that allows one
unknown coordinate. Each interface will be used twice, with different
state spaces. The related split and recovery constructions in
\cite[Section~2]{HeightThirteen,HeightFour} explain the method; the exact
statements and proofs needed here are included below.

\subsection{Monoid values and deterministic updates}

If $\Sigma=\emptyset$, its only languages are $0$ and $1$, both of
height zero. Assume henceforth that $\Sigma$ is nonempty, and fix a
morphism
\[
                         T:\Sigma^*\longrightarrow M
\]
into a finite monoid $M$. For a deterministic finite automaton, take $M$
to be the finite monoid of state transformations induced by words,
including the identity transformation. Its accepted language is a finite
union of fibers of $T$. Thus it suffices to construct a height-three
expression for each fiber of every such morphism.

Letters occupy intervals $[a,a+1)$ with integer endpoints. For an
available interval $[a,b)$, $a\le b$, write $T_{a,b}$ for its monoid
value, with $T_{a,a}=1_M$. Products are in reading order:
\[
                         T_{a,c}=T_{a,b}T_{b,c}.
\]
Let $V=\mathbf F_2^M$ have basis $(e_m)_{m\in M}$. Each $d\in M$ acts
linearly on the right by $e_m d=e_{md}$. This is the linearized right regular action
\cite[Section~3.1]{AMSV2009}. It is faithful because
$e_{1_M}d=e_d$ identifies $d$. Vectors are written on the left of their
linear maps, so $vLA$ means first apply $L$, then $A$. We will also use
$U=V\oplus\mathbf F_2$; the additional coordinate will turn affine maps
on $V$ into linear maps on $U$. All vector spaces below are finite.

Put $N_M=|M|!$. Finiteness gives
\begin{equation}\label{alg:power-periodicity}
          c^{i+N_M}=c^i\qquad(c\in M,\ i\ge |M|).
\end{equation}
Indeed two of $c^0,\ldots,c^{|M|}$ agree, say $c^a=c^b$ with
$0\le a<b\le |M|$. Powers are then periodic from exponent $a$, with
period $b-a$, which divides $N_M$. This includes $|M|=1$.

A \emph{prefix code} is a set $E\subseteq\Sigma^+$ containing no
proper-prefix pair. Its members will be called \emph{episodes}. Every
word in $E^*$ has a unique factorization into episodes: the first factors
of two proposed factorizations are comparable by the prefix relation and
therefore equal, after which induction applies. The empty word has the
empty factorization.

For $D\subseteq E$, assign to each $e\in D$ a deterministic total map
$G_e:S\to S$ on a fixed finite state set $S$. On $w=e_1\cdots e_r$,
let $G_w$ apply the updates in reading order; $G_\eps$ is the identity.
The \emph{graph language} from $x$ to $y$ on $D^*$ is
\[
                       \{w\in D^*:G_w(x)=y\}.
\]
The indicated domain is part of this definition. The update is determined
by the episode itself, independently of choices made by an expression
that tests its graph.

Finite intersections can be expressed using union and complement and do
not increase generalized star height. A finite language has height zero:
write its words as concatenations of letters, use $1$ for $\eps$, and
take a finite union. Auxiliary finite indices below are never input
letters.

\subsection{The split identity}

The next lemma separates two obligations. Every successful split must
transmit the correct state and consume exactly one episode. Existence
of a split is required only outside the domain that is to be skipped.

\begin{lemma}[Split identity]\label{alg:split}
Let $D'\subseteq D\subseteq E$, where $E$ is a prefix code, and give
each $e\in D$ a deterministic total update $G_e$ of a finite set $S$.
Suppose languages $P_x,Q_y\subseteq\Sigma^*$, $x,y\in S$, satisfy:
\begin{enumerate}[label=\textup{(\roman*)}]
\item If $w\in D^*$ has a prefix $uv$ with $u\in P_x$ and $v\in Q_y$,
then $w$ has a first episode $e$, $uv=e$, and $G_e(x)=y$.
\item For every $e\in D\setminus D'$ and every $x\in S$,
$e\in P_xQ_{G_e(x)}$.
\end{enumerate}
Let $W_{x,y}$ be the graph languages for these same updates on $(D')^*$.
Then the graph language from $x$ to $y$ on $D^*$ is
\begin{equation}\label{alg:split-formula}
\begin{split}
D^*\cap\Bigg[W_{x,y}\ \cup\ &
 \left(\bigcup_{i\in S}W_{x,i}P_i\right)
 \left(\bigcup_{j,i\in S}Q_jW_{j,i}P_i\right)^*\\[-2pt]
 &\hspace{24mm}\cdot
 \left(\bigcup_{l\in S}Q_lW_{l,y}\right)\Bigg].
\end{split}
\end{equation}
If $D$, all $P_x$, and all $Q_y$ have height at most one, while all
$W_{x,y}$ have height at most $H\ge0$, the resulting graph languages
have height at most $\max\{2,1+H\}$.
\end{lemma}

\begin{proof}
Fix a word in the right-hand side. Its membership in $D^*$ fixes the
ambient episode factorization, but does not by itself align an arbitrary
interior factor with that factorization. We establish alignment in order,
starting at the beginning of the word.

The first $W$-factor begins at a true boundary. Since its words are in
$(D')^*$, its own first episode and the ambient first episode are prefix
comparable. Both belong to the prefix code $E$, so they are equal.
Repeating this argument shows that the whole $W$-factor consumes complete
ambient episodes, and its endpoint is another true boundary. Its graph
indices describe their actual composite update.

Now consider an accepted factorization through the second term of
\eqref{alg:split-formula}. The concatenation of the first $P$-factor and
the following $Q$-factor begins at the true boundary just established.
Condition~\textup{(i)} says that their concatenation is exactly the next
episode and takes the indicated input state to the indicated output.
The next $W$-factor starts at that episode's true endpoint, so the same
prefix-code argument applies again. Alternating these two steps through
the factorization establishes every boundary and propagates every state.
Only this induction licenses alignment; no assertion about a general
interior factor in $E^*$ is needed. A zero-fold middle star leaves one
$P,Q$ pair, to which the same argument applies. The first term $W_{x,y}$
already records the desired composite on its domain.

For the reverse inclusion, follow the actual computation on a word of
$D^*$. Use condition~\textup{(ii)} to split each episode outside $D'$,
and group every intervening run of episodes from $D'$ into a $W$-factor
with its actual endpoint states. These choices form the second term in
\eqref{alg:split-formula}. If there is no episode outside $D'$, the word
belongs to $W_{x,y}$. This includes the empty word. Empty domains are
therefore allowed; if $S$ is empty there are no indexed assertions.

The height estimate can be read directly from the expression. The factor
$D^*$ has height at most two. Each factor inside the middle star has
height at most $\max\{1,H\}$, and starring it gives at most
$1+\max\{1,H\}$. Unions, concatenations, and the outer intersection
preserve the maximum of their argument heights. Thus the bound is
$\max\{2,1+H\}$.
\end{proof}

When $D'=\emptyset$, the skip language $W_{x,y}$ is $1$ for $x=y$
and $0$ otherwise. Its height is zero, so this case of the lemma has
height bound two. Notice also that condition~\textup{(i)} rules out a
successful pair on the empty word. Independent choices in $P_x$ and
$Q_y$ must satisfy it even when their proposed metadata disagree.

\subsection{Removing affine corrections}

We will deliberately alter an episode's translation term to make its
split possible. The following identity explains why this is harmless.
The output on zero isolates the translation of a whole composite, even
when none of its linear factors is invertible. The two computations must
be required on the same word, which is why the expression uses an
intersection.

\begin{lemma}[Boolean affine recovery]\label{alg:affine-recovery}
Let $S$ be a finite vector space and let the episode updates on
$D\subseteq E$ be
\[
                          G_e(v)=vL_e+b_e,
\]
with $L_e\in\operatorname{End}_{\rm lin}(S)$ and $b_e\in S$ determined
by $e$. If $R_{x,z}$ are the graph languages of their composites on
$D^*$, then the graph language for the product of their linear parts,
from $x$ to $y$, is
\begin{equation}\label{alg:linear-recovery-formula}
                  \bigcup_{b\in S}
                         \bigl(R_{x,y+b}\cap R_{0,b}\bigr).
\end{equation}
A common height bound for the $R_{x,z}$ is preserved.
\end{lemma}

\begin{proof}
Composition in reading order obeys
$(vL+b)A+c=vLA+(bA+c)$. Induction, starting from the identity on the
empty word, therefore writes each word's composite as $v\mapsto vL+b$,
where $L$ is precisely the product of its episode linear parts.
For that word, the output at zero is $b$, and the output at $x$ is
$y+b$ exactly when $xL=y$. Requiring both facts on the word gives the
intersection in \eqref{alg:linear-recovery-formula}; taking the union
allows every possible translation $b$. These are finite Boolean
operations, so they preserve the asserted height bound.
\end{proof}

\begin{lemma}[Homogeneous lift]\label{alg:homogeneous-lift}
For a finite vector space $S$ over a finite field $\mathbf F$, the map
$F(v)=vL+b$ has the linear lift on $S\oplus\mathbf F$
\begin{equation}\label{alg:lift-formula}
                     \widehat F(v,c)=(vL+cb,c).
\end{equation}
Products of these lifts are the lifts of the corresponding composites
in reading order. In particular, graph languages for products of the
lifts give those of the affine composites by using inputs $(v,1)$
and outputs $(y,1)$, without increasing height.
\end{lemma}

\begin{proof}
The displayed map is linear in $(v,c)$. For
$F(v)=vL+b$ followed by $G(v)=vA+d$, the two lifted maps give
\[
            (v,c)\longmapsto(vLA+c(bA+d),c),
\]
which is the lift of their affine composite. Induction proves the
product assertion; the slice $c=1$ proves the graph assertion. The
empty product is the identity in both spaces. No invertibility is
required. We will use this with $S=V$, $\mathbf F=\mathbf F_2$, and
$S\oplus\mathbf F=U$.
\end{proof}

\subsection{A shift table that hides one entry}

A cut can leave one interval product unreadable by either factor. We
will replace the demand to know that product by a finite certificate that
checks all of its possible values. To make such a certificate available,
we need a deterministic affine correction with a constant coordinate
line through every tuple, for every input state.

The order of quantifiers matters: the length is selected first, uniformly
in the table of linear maps; a correction is then fixed for each table;
only afterward may a tuple and input state select a successful coordinate.
The product-action version in \cite[Lemma~2.3]{HeightFour} motivates this
step, but an arbitrary table is needed for the second split below.

\begin{lemma}[Shift table]\label{alg:shift-table}
For every nonempty finite set $A$ and every finite vector space $S$
there is an integer $l\ge1$, depending only on $|A|$ and $|S|$, with
the following property. For every function
\[
                    f:A^l\longrightarrow
                                    \operatorname{End}_{\rm lin}(S)
\]
there is a function $\beta_f:A^l\to S$ such that, for every
$\mathbf a=(a_1,\ldots,a_l)\in A^l$ and $x\in S$, there is an
$i\in\{1,\ldots,l\}$ for which
\begin{equation}\label{alg:shift-constancy}
 b\longmapsto x f(a_1,\ldots,a_{i-1},b,a_{i+1},\ldots,a_l)
       +\beta_f(a_1,\ldots,a_{i-1},b,a_{i+1},\ldots,a_l)
\end{equation}
is constant on $A$. The integer $l$ is uniform over $f$, and
$\beta_f$ is independent of $x$.
\end{lemma}

\begin{proof}
Put $a=|A|$ and $q=|S|$. If $a=1$ or $q=1$, take $l=1$ and
$\beta_f=0$; every required function is constant. Assume $a,q>1$.
For an integer $l$ to be chosen, choose independent uniform vectors
$\beta(\mathbf c)\in S$ at all cells $\mathbf c\in A^l$. Fix any
table $f$. For each $x\in S$ and $\mathbf a\in A^l$, let
$\mathcal B_{x,\mathbf a}$ be the event that none of the $l$
coordinate lines through $\mathbf a$ has the asserted constancy.

Condition on the value of $\beta(\mathbf a)$. For a specified line,
constancy forces exactly one value of $\beta$ at each of its other
$a-1$ cells, regardless of $f$. The line therefore succeeds with
probability $q^{1-a}$. Different lines share only their central cell,
so these successes are conditionally independent. The conditional
failure probability is independent of the central value, giving
\[
                    \Pr(\mathcal B_{x,\mathbf a})
                          =(1-q^{1-a})^l=:p_l.
\]

This event depends only on random values at Hamming distance at most
one from $\mathbf a$, where Hamming distance counts unequal entries.
Events with centers at distance greater than two use disjoint random
variables. Hence each event is independent of the joint family of all
such distant events. A dependency neighborhood has size at most
\begin{equation}\label{alg:dependency-bound}
 q\left(1+l(a-1)+\binom l2(a-1)^2\right)-1
              \ \le\ D_l:=q(1+la+l^2a^2).
\end{equation}
The factor $q$ includes every possible input vector at each center.

We use the following form of the finite dependency argument of
Erd\H{o}s and Lov\'asz~\cite[Section~2]{ErdosLovasz1975} and include its proof. Suppose each event of a finite family
is independent of the joint family of its nonneighbors, has at most
$D$ neighbors, and has probability at most $s(1-s)^D$, where
$0<s<1$. Then the probability that all the events fail to occur is
positive. Here is the conditional-probability argument, to make the
precise independence requirement explicit.

Induct on the size of a set of other events being conditioned not to
occur. Along with positivity of these conditioning events, prove that
the conditional probability of any remaining event is at most $s$.
For an empty set this follows from the assumed bound. Positivity for
a set of size $k$ follows by deleting one event and using the bound
for a conditioning set of size $k-1$. To prove the conditional bound
at size $k$, split the conditioning set into the neighbors $J_1$ and
nonneighbors $J_0$ of the event $B$ under consideration. If $J_1$ is
empty, joint independence gives its unconditional probability. Otherwise,
writing $C_i$ for simultaneous nonoccurrence of $J_i$, we have
\[
 \Pr(B\mid C_0\cap C_1)
 \le \frac{\Pr(B\mid C_0)}{\Pr(C_1\mid C_0)}
 \le \frac{\Pr(B)}{(1-s)^{|J_1|}}
 \le s.
\]
For the denominator, expose the events in $J_1$ one at a time; every
conditioning set then has size at most $k-1$, so the induction bounds
each nonoccurrence probability below by $1-s$. The numerator uses
joint independence from $J_0$. This completes the induction, and
successive conditioning gives positive probability of full avoidance.

Apply this estimate with $D=D_l$ and
$s=1/(2(D_l+1))$. Bernoulli's inequality gives
\[
 (1-s)^{D_l}\ge1-D_ls\ge\tfrac12,
 \qquad s(1-s)^{D_l}\ge\frac1{4(D_l+1)}.
\]
The probability $p_l$ decays exponentially in $l$, whereas $D_l$
grows quadratically. Choose $l$ so that
$p_l\le1/(4(D_l+1))$. This choice depends only on $a,q$, and works
for every table $f$. Avoidance of all $\mathcal B_{x,\mathbf a}$
supplies the required $\beta_f$ simultaneously for all cells and
input vectors.
\end{proof}

Fix one successful $\beta_f$ for each of the finitely many possible
functions $f$ at these dimensions. The resulting map
$x\mapsto xf(\mathbf a)+\beta_f(\mathbf a)$ is now determined by the
tuple and table, before any splitting choices are made. A prefix can
certify its output using the other entries by checking every replacement
for the omitted coordinate. The lemma promises a usable coordinate for
each actual input, not one coordinate usable for all inputs at once.
Nor does it restrict $f$ to tables obtained from feasible word prefixes.
This unrestricted table domain will be essential when a suffix computes
hypothetical prefix values in Section~\ref{inner:section}.
